\cotefile{135}
\scanpage{17}
% transcription
pour \unsure{deviner|décrire|choisir} $P$ (Pic. d'invariants $L,N$) telle
que $P\mid\mathbb{R}$ soit splittée -- ou ce qui revient au
$=$, telle que $P$ provienne d'une cat. de Pic.
d'inv. $M,N$, \deleted{\unsure{D'où peut on en}} \inserted{On vérifie immédiatement}
que $w(\sigma)$ ne dépend pas du choix particulier
de $L\to M\mathbin{-}0$, et que sa formation est
\inserted{fonctorielle en $M,N$ et est} compatible avec \deleted{\unsure{celle des}} images inverses
par morphismes de topos. Or se plaçant
dans la situation universelle de $M\to{}_{2}N$,
\deleted{\unsure{on trouve}} soit $\mathcal{M}\mathrel{\overset{\sigma}{-}}{}_{2}\mathcal{N}$ \inserted{sur $B$}, on trouve
aisément $\operatorname{Ext}^{i}(B;\mathcal{M},\mathcal{N})=0$ pour $i>0$,
en particulier $w(\sigma)\in\operatorname{Ext}^{3}(B;\mathcal{M},\mathcal{N})\deleted{=0}$
est nul, donc $w(\sigma)$ est nul dans tous les
cas !
% archive sheet number: 16

\scanpage{18}
% transcription
\mspage{6}
$M,N$ faisceaux abéliens sur topos $X$

$\mathscr{P}$ 2-catégorie des chps de Picard sur $X$
épinglés par $P$. $P$ est muni de
la composition de Baer.

$\pi_0(P)=$ groupe des objets de $P$ \unsure{à} équivalence
près
\[
0\to\operatorname{Ext}^{2}(M,N)\to\pi_0(P)\to\operatorname{Hom}(M,{}_{2}N)\to0
\]

$\pi_1(P)=$ groupe des classes d'isomorphie
d'autoéquivalences d'un objet de
$\mathscr{P}$
\[
=\operatorname{Ext}^{1}(M,N)
\]

$\pi_0(P)=$ groupe des automorphismes du
foncteur identique d'un
objet $P$ de $\mathscr{P}$
\[
=\operatorname{Hom}(M,N)
\]

$\mathscr{P}$ 2-catégorie de Picard \uline{stricts} ?

$P\otimes P=$ catégorie dont les objets sont
$(L,L')$ avec \unsure{et|où|$\partial$} $L\simeq L'$ \uline{localement}
\[
\operatorname{Isom}\big((L,L'),(U,U')\big)=
\frac{\operatorname{Isom}(L,U)\times\operatorname{\unsure{Ism|Isom|Hom}}(L',U')}{\underset{\swarrow\,(\text{diagonale})}{N}}
\]
% archive sheet number: 17

\scanpage{19}
% transcription
\mspage{7}
\diagram{\resizebox{\linewidth}{!}{\begin{tikzcd}[ampersand replacement=\&,column sep=small,row sep=small]
{\left.\begin{gathered}\text{la 3-catégorie des}\\\text{Gr-catégories }C,\\\text{avec }\pi_0(C)=G\\\pi_1(C)=\deleted{\unsure{Z}}M\\[2pt]\text{et}\\\text{2-catégorie des}\\\text{splittings}\\\text{des Gr-catégories}\end{gathered}\right\}}
\arrow[rr,leftrightarrow,"{?}"] \arrow[rr,shift right=4,"{\deleted{\text{``noyaux''}}}"',"\deleted{\substack{\text{de }G\text{-liens}\\\text{de centre }Z}}" description]
\& {} \&
{\begin{gathered}\text{la 3-catégorie des}\\\text{2-gerbes sur }B_G\\\deleted{\text{\unsure{de faisceaux|de liens}}}\\\deleted{\text{\illegible{2}}}\\\deleted{\text{\unsure{de lien|de liens}}}\\\deleted{\text{\unsure{de centre }}Z}\\\text{de lien commutatif}\\\text{donné }(\deleted{\unsure{Z}}M\text{ avec }G\\\text{opérant dessus})\end{gathered}}\\
{} \& {\begin{gathered}\text{``noyaux'' et}\\\text{représentations}\\\text{des noyaux}\\\text{par extensions}\end{gathered}}\arrow[ul]\arrow[r,leftrightarrow,very thick]\&
{\begin{gathered}\text{Liens sur } B_G\\\text{et 2-gerbes}\\\text{des représentations}\\\text{des Liens}\end{gathered}}\arrow[u]\arrow[l,shift left=3]
\end{tikzcd}}}

Soit $C$ une Gr-catégorie d'invariants $G,M$. On va lui associer
une 2-gerbe \inserted{$\widetilde C$} sur $B_G$, donc à tt objet \inserted{$X$} de $B_G$
(un $G$-ens) on va associer une 2-catégorie. Si $X=$ objet
final de $B_G$, ce sera la 2-catégorie des \uline{splittings} de
$C$.

\deleted{Fixons pour tt $x\in G$ un représentant $L(x)\in\operatorname{Ob} C$,}
\deleted{et pour} \deleted{Un $L$-splitting de $C$ est la donnée pour
tout couple $(x,y)$, d'un isom $\varphi_{x,y}:L(x)\otimes L(y)\simeq L(xy)$,} \deleted{Un}
\deleted{\unsure{trivialisation du couple|modification du couple}} \deleted{$(L,\varphi)$
$(L:G\mathbin{\unsure{-|\to}}\operatorname{Ob}C,\ \varphi:G\timesG\mathbin{\unsure{-|\to}}\mathrm{Fl}(C)$,}
\deleted{est \unsure{des} choix d'une 2-cochaîne \unsure{de} de telle façon que}
\deleted{$\lambda:G\times G\longrightarrow M$ la 3-\unsure{cochaîne|cochaîne d'obstruction} $h$ qui exprime}
\deleted{telle que $\unsure{\psi_{x,y}=}$ le système des \inserted{En termes de $L,\varphi$}}
\deleted{$\bigl[\begin{smallmatrix}L:G\,-\,\operatorname{Ob}C\\\varphi:G\times G\,-\,\mathrm{Fl}C\end{smallmatrix}\bigr]$}
\deleted{$\psi_{x,y}=\varphi_{x,y}\lambda(x,y)$ le défaut d'associativité de}
\deleted{\unsure{foncteur (avec $L$)}} \deleted{$C$, soit \unsure{une} \uline{nulle}. L'ens des $L$}
\deleted{splittings forme un ens. principal homogène sous}
\deleted{$Z^2(G,M)$. Soient $(L,\varphi)$, $(L',\varphi')$ deux splittings}

Soit $C_0$ \deleted{une} \inserted{la} (\deleted{Gr})-catégorie \uline{réduite} \uline{splittée}
\deleted{(ceci revient à dire l'isomorphisme d'associativité est tjrs l'identité)}
\inserted{épinglée de type $G,M$,}
\[
\begin{aligned}
&\operatorname{Ob}C=G\quad\mathrm{Fl}\,C=G\times M,\quad \deleted{\unsure{\mathrm{Hom}(g,g')}}\ G\times M\underset b{\overset s{\rightrightarrows}}G\quad s=b=\mathrm{pr}_1\\
&(g,m)\circ(g,m')=(g,m+m'),\quad\deleted{(g,m)\otimes(g',m')}\ g\otimes g'=gg'\\
&(g,m)\otimes(g',m')=(gg',m+m'),\quad a_{g,g',g''}=\mathrm{id}_{gg'g''}.
\end{aligned}
\]
On appellera \uline{splittings} de $C$ toute équivalence de Gr-catégories
\inserted{épinglées} $S\inserted{=(S_0,\varphi)}:C_0\longrightarrow C$, \deleted{\unsure{fon}} (donc donné par un
\inserted{\unsure{fon}}
% archive sheet number: 18

\scanpage{20}
% transcription
\deleted{$(\unsure{L},\varphi_{x,y})$ \unsure{et les}}

\deleted{La catég} Les splittings de $C$ forment une catégorie,
qui est un 1-torseur sous Gr-catégorie des auto-équivalences de la Gr-catégorie $C_0$ avec elle $=$.
Quels sont les invariants de cette dernière
Gr-catégorie ?

\uline{Objets} Les applications $\varphi:G\times G\mathbin{-}M$ (2-cochaines)
qui sont des \unsure{cbords|cobords}

\uline{Morphismes} $\varphi\to\psi$ est l'ens des 1-cochaines
dont le cobord est $\psi-\varphi$

Objets à isom près \quad $H^2(G,M)$

endomorphismes d'un objet : $Z^1(G,M)$ \deleted{$\unsure{=}$}

\diagram{\begin{tikzcd}[ampersand replacement=\&,column sep=large]
C_0 \arrow[r,shift left=1.7,"S",""{name=up,below}] \arrow[r,shift right=1.7,"S'"',""{name=lo,above}] \& C
\arrow[from=up,to=lo,Rightarrow,shift left=2,"\alpha'" description] \arrow[from=up,to=lo,Rightarrow,shift right=2,"\alpha"']
\end{tikzcd}}
\diagram{\begin{tikzcd}[ampersand replacement=\&,column sep=small]
\deleted{\unsure{\varphi}}\arrow[d] \& L_x,\varphi_{x,y}\arrow[d,shift right=1,"\alpha_x"']\arrow[d,shift left=1,"\alpha'_x"]\& \deleted{\unsure{\varphi\in\mathrm{Fl}\,C_M}}\\
{} \& L_{x'},\varphi_{x',y'} \& {}
\end{tikzcd}}
\[
\begin{aligned}
&\deleted{\unsure{\alpha,2}}\qquad\deleted{\alpha\to\alpha'}\quad\alpha'(x)=\alpha(x)\varphi(x)\\
&\hspace{12em}\searrow\ \delta\varphi=0\\
&\operatorname{Hom}(\alpha,\alpha')=\bigl\{\deleted{\unsure{g\in G}}\inserted{m\in M}\mid\deleted{\unsure{g=\alpha'\alpha^{-1}}}\\
&\hspace{13em}\deleted{\alpha'=\alpha(\delta g)}\\
&\hspace{10em}\alpha'(x)=\alpha(x)(gm-m)
\end{aligned}
\]

\noindent\rule{.45\linewidth}{.4pt}

Soit $C$ une \inserted{(Gr-)-}catégorie \deleted{\unsure{de}} \inserted{d'}\unsure{invar|invariants} $=G,M$
$X$ un ens. où $G$ opère. On va construire une
2-catégorie $\unsure{S|\mathscr{S}}(C,X)$ :

0-objets : Systèmes $L,\varphi$ \unsure{associant|consistant|ainsi}

1) à tt couple $(x,g)$ $(x\in X,g\in G)$ un $L(x,g)\in\operatorname{Ob}{C_g}$

2) à tt triple $(x,g,g')$ $(x\in X,g\in G,g'\in \unsure{G'|G})$ un isom
\[
L(x,g)\otimes L(\deleted{\unsure{x,g'}}\inserted{gx,g'})\simeq L(x,g'g)
\]
\margin{\diagram{\resizebox{\linewidth}{!}{\begin{tikzcd}[ampersand replacement=\&,row sep=small]X\times G\arrow[d,shift right=1,"\mathrm{pr}_1"']\arrow[d,shift left=1,"\pi"]\\X\end{tikzcd}}}}
\margin{\diagram{\resizebox{\linewidth}{!}{\begin{tikzcd}[ampersand replacement=\&,column sep=tiny]x\arrow[r]\&g.x\arrow[r]\&g'gx\end{tikzcd}}}}
% archive sheet number: 19

\scanpage{21}
% transcription
\mspage{1}\mspage{8}
$\Gamma$ catégorie

$M$ groupe \deleted{\unsure{commutatif}}

Une \uline{extension} de $\Gamma$ par $M$ est la donnée

a) Pour tte flèche $u\in\mathrm{Fl}\,\Gamma$, d'un $M$-\inserted{bi}torseur $P_u$

b) Pour deux flèches composables $(u,v)$ de $\Gamma$, d'un isom
\[
\varphi_{u,v}:P_u\otimes P_v\simeq P_{uv}
\]
\margin{\diagram{\begin{tikzcd}[ampersand replacement=\&,row sep=small]
\mathrm{Fl}_3(\Gamma)\arrow[d,shift left=3]\arrow[d,shift left=1]\arrow[d,shift right=1]\arrow[d,shift right=3]\\
\mathrm{Fl}_2(\Gamma)\arrow[d,shift left=2]\arrow[d]\arrow[d,shift right=2]\\
\mathrm{Fl}_1\Gamma\arrow[d,shift left=1]\arrow[d,shift right=1]\\
\mathrm{Fl}_0\Gamma
\end{tikzcd}}}
avec axiome de commutativité du diagramme
suivant, si $(u,v,w)$ composables (axiome \unsure{pentago.|pentagon.|pentagonal}?)
\diagram{\begin{tikzcd}[ampersand replacement=\&]
(P_u\otimes P_v)\otimes P_w\arrow[r,"{\varphi_{u,v}\otimes\mathrm{id}}","\sim" description]\arrow[d,"\mathrm{can}"] \& P_{uv}\otimes P_w\arrow[d,"{\varphi_{uv,w}}"]\\
P_u\otimes(P_v\otimes P_w)\arrow[d,"{\mathrm{id}\otimes\varphi_{v,w}}"] \& P_{uvw}\\
P_u\otimes P_{vw}\arrow[ur,"{\varphi_{u,vw}}"'] \& {}
\end{tikzcd}}

[ \uline{NB} Au lieu de se donner $M$, il suffit de se donner
une \inserted{Gr-}catégorie \deleted{de Picard} \inserted{$\mathscr{P}$} -- ici la catégorie des
$M$-bitorseurs ]

[ \uline{Invariants} \textcircled{1}
\diagram{\begin{tikzcd}[ampersand replacement=\&]
\mathrm{Fl}(\Gamma)\arrow[r,"\alpha"]\arrow[dr,bend right=15,"\alpha'"']\& \pi_0(\mathscr{P})\quad(=\operatorname{Aut\,ext}(M))\\
{} \& \deleted{\unsure{\operatorname{Aut}(\mathscr P)}}\ \underset{\operatorname{Aut}(M)}{\text{\unsure{'' $\sim$ si $M$ comm|= si $M$ commutatif}}}
\end{tikzcd}}
\textcircled{2}
\[
c\in\underbrace{H^2(\Gamma;\deleted{\unsure{Z}}\mathscr P;\deleted{\unsure{\alpha}}\alpha)}_{
\left\{\begin{gathered}\text{c'est un }H^{2}\text{ non}\\
\text{commutatif, en général}\\
\text{sans élément neutre }\deleted{\unsure{?}}\end{gathered}\right.}
\]
]

\noindent\rule{.4\linewidth}{.4pt}

[ Une 2-extension de $\Gamma$ par $M$ \inserted{[\uline{commutatif}]} est le choix :

a) Pour tte flèche $u\in\mathrm{Fl}\,\Gamma$, d'une $M$-gerbe $P_u$

b) Pour deux flèches composables $u.v$, d'une
équivalence
\[
\varphi_{u,v}:P_u\otimes P_v\simeq P_{uv}
\]

c) Pour trois flèches composables, d'un isom
rendant commutatif le diagramme
pentagonal
% archive sheet number: 20

\scanpage{22}
% transcription
\[
\deleted{P_u\otimes P_v\otimes P_w\longrightarrow P_{uvw}}
\]
\[
\begin{aligned}
&\deleted{P\otimes Q\otimes R\mathbin{\unsure{-|\to}}S}\\
&\left\{\begin{aligned}
&\deleted{P\otimes Q-S\otimes R^{-1}}\\
&\deleted{Q\otimes R-P^{-1}\otimes S}\\
&\deleted{R\otimes S^{-1}-Q^{-1}\otimes P^{-1}}\\
&\deleted{S^{-1}\otimes P-R^{-1}\otimes Q^{-1}}
\end{aligned}\right.
\end{aligned}
\]
\diagram{\begin{tikzcd}[ampersand replacement=\&,column sep=small]
{}\arrow[r,"{\deleted{w}}"]\&{}\arrow[r,"{\deleted{v}}"]\&{}\arrow[r,dash,"{\deleted{w}}"]\&{}
\end{tikzcd}}
\[
\deleted{L_u\otimes L_v\otimes L_w\simeq L_{uvw}}
\]
\diagram{\begin{tikzcd}[ampersand replacement=\&,row sep=small]
\deleted{u,v,w}\arrow[d]\\\deleted{u,v}
\end{tikzcd}}
\[
\deleted{L_u\otimes L_v-L_{uv}}
\]
\diagram{\begin{tikzcd}[ampersand replacement=\&,row sep=small,column sep=small]
{}\&\deleted{\mathrm{Fl}_3\Gamma}\arrow[d,dash,shift left=3]\arrow[d,dash,shift left=1]\arrow[d,dash,shift right=1]\arrow[d,dash,shift right=3]\&{}\\
\deleted{M_{\Gamma_2}}\arrow[r,dash]\arrow[d,dash,shift left=2]\arrow[d,dash]\arrow[d,dash,shift right=2]\&\deleted{\mathrm{Fl}_2(\Gamma)}\arrow[r,equals]\arrow[d,dash,shift left=2]\arrow[d,dash]\arrow[d,dash,shift right=2]\&\deleted{\mathrm{Fl}\,\Gamma\times_{\unsure{p|P|\mathrm{Ob}}}\mathrm{Fl}\,\Gamma}\arrow[d,dash,shift left=2]\arrow[d,dash]\arrow[d,dash,shift right=2]\\
\deleted{M_{\Gamma_1}}\arrow[r]\arrow[d,shift left=1]\arrow[d,shift right=1]\&\deleted{\mathrm{Fl}_1(\Gamma)}\arrow[r,equals]\arrow[d,shift left=1]\arrow[d,shift right=1]\&\deleted{\mathrm{Fl}\,\Gamma}\arrow[d,dash,shift left=1]\arrow[d,dash,shift right=1]\\
\deleted{M_{\Gamma_0}}\arrow[r]\&\deleted{\mathrm{Fl}_0(\Gamma)}\arrow[r,equals]\&\deleted{\operatorname{Ob}\Gamma}
\end{tikzcd}}
\diagram{\resizebox{\linewidth}{!}{\begin{tikzcd}[ampersand replacement=\&,column sep=small,row sep=small]
\deleted{\Gamma_3}\arrow[r]\arrow[d,dash,shift left=3]\arrow[d,dash,shift left=1]\arrow[d,dash,shift right=1]\arrow[d,dash,shift right=3]\&
\deleted{\mathscr{P}_3}\arrow[d,dash,shift left=3]\arrow[d,dash,shift left=1]\arrow[d,dash,shift right=1]\arrow[d,dash,shift right=3]\& \deleted{\substack{=\text{\unsure{accouplement|accouplements} de trois objets}\\\text{vers un quatrième}}}\\
\deleted{\Gamma_2}\arrow[r]\arrow[d,dash,shift left=2]\arrow[d,dash]\arrow[d,dash,shift right=2]\&\deleted{\mathscr{P}_2}\arrow[d,dash,shift left=2]\arrow[d,dash]\arrow[d,dash,shift right=2]\&\deleted{\substack{=\text{\unsure{accouplement|accouplements} de deux objets}\\\text{de }P\text{ vers un troisième}}}\\
\deleted{\Gamma_1}\arrow[r]\arrow[d,dash,shift left=1]\arrow[d,dash,shift right=1]\&\deleted{\mathscr{P}_1}\arrow[d,dash,shift left=1]\arrow[d,dash,shift right=1]\&\deleted{=\operatorname{Ob}P}\\
\deleted{\Gamma_0}\arrow[r]\&\deleted{e}\&{}
\end{tikzcd}}}
\[
\begin{aligned}
&\deleted{H^2(\Gamma,M)}\\&\deleted{H^1(\Gamma,M)}\\&\deleted{H^0(\Gamma,M)}
\end{aligned}
\]
\deleted{$\mathscr P_I=$ ensemble des} \deleted{couples} \inserted{\deleted{triples}} \deleted{$(L,M,\varphi)$}
\[
\begin{aligned}
&\deleted{I\hookrightarrow\text{ens fini tot. ordonné}}\\
&\deleted{L:I\to\operatorname{Ob}P\quad L=(L_i)_{i\in I}}\\
&\deleted{M\in\operatorname{Ob}\unsure{P|\mathscr P}}\\
&\deleted{\varphi:\bigotimes_{i\in I}L_i-M}
\end{aligned}
\]
\[
\begin{aligned}
&\deleted{H^2}\\
&\deleted{L_u L_v\simeq L_w}\\
&\deleted{\unsure{L_{v}}}\qquad\deleted{L_1\otimes L_2-M}\qquad\deleted{\Gamma_1\to\pi_0(P)}\\
&\deleted{L_1\otimes L_2\otimes M^{-1}\simeq}\\
&\deleted{L_1\otimes\cdots\otimes L_n\otimes S^{-1}\simeq1}
\end{aligned}
\]
\diagram{\resizebox{\linewidth}{!}{\begin{tikzcd}[ampersand replacement=\&,column sep=small,row sep=small]
{}\&{}\arrow[dr,leftrightarrow,"{\deleted{v}}"]\&{}\\
{}\arrow[ur,"{\deleted{w}}"]\arrow[rr,"{\deleted{u}}","{\deleted{L_u}}"']\&{}\&{}
\end{tikzcd}}}
\diagram{\begin{tikzcd}[ampersand replacement=\&,column sep=small,row sep=small]
{}\&{}\&{}\\
{}\&{}\&{}\arrow[ul,"{\deleted{w}}"']\\
{}\arrow[uur,"{\deleted{w'}}"]\arrow[urr,"{\deleted{v'}}",near end]\arrow[rr,"{\deleted{u}}"',"{\deleted{L_u}}"',near end]\&{}\&{}\arrow[u,"{\deleted{v}}"',"{\deleted{L_v}}"]\arrow[uul,"{\deleted{u'}}"']
\end{tikzcd}}
\[
\begin{aligned}
&\deleted{P\otimes Q\to R}\qquad \deleted{\unsure{P}-Q\otimes R}\\
&\deleted{R\otimes Q^{-1}\to P}\qquad\deleted{\unsure{Q}-\mathord{\cdot}}\\
&\deleted{{\unsure{p|P}}^{-1}\otimes R\to Q}\\
&\deleted{L_u L_v\sim L_{v'}}\qquad\deleted{L_v L_w\sim L_{u'}}\\
&\deleted{L_{\unsure{v'|w'|uv}}L_w\sim L_{w'}}\qquad\deleted{L_u L_{u'}-L_{w'}}\\
&\deleted{L_u L_v L_w\sim L_{w'}}\qquad\deleted{L_u L_v L_w-L_{w'}}
\end{aligned}
\]
% archive sheet number: 21

\scanpage{23}
% transcription
\mspage{2}\mspage{9}
\diagram{\begin{tikzcd}[ampersand replacement=\&]
(P_u\otimes P_v)\otimes P_w\arrow[r,"{\varphi_{u,v}\otimes\mathrm{id}}","\sim" description]\arrow[d,leftrightarrow,"\mathrm{can}"]\&P_{uv}\otimes P_w\arrow[d,"{\varphi_{uv,w}}",""{name=aa,at start}]\\
P_u\otimes(P_v\otimes P_w)\arrow[d,"{\mathrm{id}\otimes\varphi_{v,w}}"]\&P_{uvw}\\
P_u\otimes P_{vw}\arrow[ur,"{\varphi_{u,vw}}"',""{name=bb,near end}]\&{}
\arrow[from=aa,to=bb,Rightarrow,bend right=18,shorten >=2pt,"{\alpha_{u,v,w}}"']
\end{tikzcd}}

d) Un axiome de compatibilité quand
on a \uline{quatre} flèches composables

\deleted{$\unsure{uvwz}\quad\unsure{tuvwz}$}
\diagram{\begin{tikzcd}[ampersand replacement=\&,row sep=small]
\deleted{((P_t\otimes P_u)\otimes P_v)\otimes P_{\unsure{z|w}}}\arrow[d,"{\deleted{\mathrm{can}}}"]\\
\deleted{(P_t\otimes(P_u\otimes P_v))\otimes P_w}\arrow[d,"{\deleted{\unsure{\exists\,\mathrm{can}|\mathrm{can}}}}"]\\
\deleted{P_t\otimes(P}
\end{tikzcd}}

On veut en fait que tous les isomorphismes
entre foncteurs de la
forme
\diagram{\begin{tikzcd}[ampersand replacement=\&]
(\cdots((P_{u_1}\otimes P_{u_2})\otimes P_{u_3})\cdots P_{u_n})\arrow[dr]\&{}\\
{}\&P_{u_1u_2\cdots u_n}
\end{tikzcd}}
\[
\deleted{P_t\otimes(P_u\otimes(P_v\otimes P_w))}
\]
déduits des $\varphi_{u,v}$ \deleted{\unsure{et des}} \inserted{et des}
isom d'associativité \inserted{canoniques}
forment un système transitif \deleted{\unsure{du}}
\deleted{\unsure{quand tt}}

\uline{NB} Au lieu de $M$, on peut aussi se donner
une 2-\deleted{Gr-}cat \inserted{de Picard} $\mathscr P$ (p. ex. celle des $M$-gerbes).
\[
\begin{aligned}
&\deleted{\mathrm{Fl}_1(\Gamma)\mathbin{\unsure{-|\to}}\pi_0(\mathscr P)}\qquad
\deleted{\text{(\unsure{zéro|2o} en l'occurrence)}}\\
&\deleted{\mathrm{Fl}_2(\Gamma)\to}
\end{aligned}
\]
]
% archive sheet number: 22

\scanpage{24}
% transcription
\mspage{3}
$\Gamma$ cat. \quad $P$ \unsure{Gr-Cat|Gr-cat.}

\diagram{\begin{tikzcd}[ampersand replacement=\&]
\mathrm{Fl}(\Gamma)\arrow[r,"\alpha"]\arrow[rr,bend right=15,"{\deleted{\unsure{\bar\alpha}}\ \bar\alpha}"']\&\pi_0(P)\arrow[r,dash]\&\operatorname{Aut}(\pi_1(P))
\end{tikzcd}}
On va définir la \uline{2-catégorie} des $\alpha$-extensions
de $\Gamma$ par $P$

\textcircled{1} 0-objets : les ext. de $\Gamma$ par $P$ subordonnées à $\alpha$

\textcircled{2} 1-objets \quad $\{P_u,\varphi_{u,v}\}\longrightarrow\{P'_u,\varphi'_{u,v}\}$
sont les systèmes d'hom
\[
\lambda(u):P_u\to P'_u
\]
donnant lieu à diagrammes commutatifs
\deleted{avec les pentagones $\Pi_{u,v,w}$ et $\Pi'_{u,v,w}$}
\inserted{$\varphi_{u,v},\varphi'_{\unsure{uv|u,v}}$}
\diagram{\begin{tikzcd}[ampersand replacement=\&]
P_u\otimes P_v\arrow[r,dash,"\sim","{\varphi_{u,v}}"']\arrow[d,"{\lambda(u)\otimes\lambda(v)}"']\&P_{uv}\arrow[d,"{\lambda(uv)}"]\\
P'_u\otimes P'_v\arrow[r,dash,"{\varphi'_{u,v}}"']\&P'_{uv}
\end{tikzcd}}

\uline{NB} Si on a un tel $\{\lambda(u)\}$, les
autres \deleted{$\lambda'(u)$} $\{\mu(u)\}$ sont de
la forme $\lambda(u)\otimes f(u)$, $f:\mathrm{Fl}\,\Gamma\to\pi_1(P)$,
avec $f(uv)=f(u).(\alpha(u).f(v))$, i.e. $f$ étant
une 1-cochaine sur $\Gamma$ à valeurs dans
$(\pi_1(P),\bar\alpha)$. Ainsi l'ens des hom
de $\{P_u,\varphi_{u,v}\}$ dans $\bigl\{P'_u,\varphi'_{u,v}\bigr\}$ est un
\deleted{\unsure{princip}} pseudo-torseur sous $Z^1(\Gamma,(\pi_1(P),\bar\alpha))$.

\textcircled{3} 2-objets \quad Soit $(\lambda_u),(\mu_u)$ \deleted{$\unsure{\leftrightarrow|\to}$} :
$\{P_u,\varphi_{u,v}\}\mathrel{\substack{\longrightarrow\\[-4pt]\longrightarrow}}\{P'_u,\varphi'_{u,v}\}$

Donc $\mu_{\unsure{u|x}}=\lambda_u\otimes f(u)$ \quad $f(u)\in Z^1(\Gamma,\pi_1(P))$.

On pose que l'ens des flèches de
\deleted{$\unsure{f}$} $\lambda$ dans $\mu$ est l'ens des $m\deleted{\unsure{f}}\in Z^0(\Gamma,\pi_1(P))$
tels que $f=\delta m$ i.e.
\diagram{\begin{tikzcd}[ampersand replacement=\&]x\arrow[r,"u"]\&y\end{tikzcd}}\hfill\textquotedblright
\[
\mu_u=\lambda_u\otimes(m(y)-\alpha(u).m(x))\qquad(?)
\]
% archive sheet number: 23

\scanpage{26}
% transcription
1. $n$ catégories (strictes -- pas strictes) et $n$-foncteurs (stricts ou pas).
\inserted{Il y a des chances que les ``pas stricts'' se ramènent aux stricts}
les $n$-catégories forment une $(n+1)$-catégorie, les $\underline{\operatorname{Hom}}$ y sont des
$n$-catégories ; notion de $i$-fidélité $(0\leq i\leq n+1)$ -- la $(n+1)$-fidélité
s'appelle $n$-équivalence.

Relations entre $n$-catégories et $(n+1)$-catégories \inserted{réduites} (au moins \unsure{en|dans les} 2 sens)
Une Gr-catégorie ``est'' une 2-catégorie ayant un seul 0-objet.

2. Exemples liés aux structures semi-simpliciales.
$\left\{\begin{gathered}\text{Espaces topologiques, groupes top, anneaux top etc}\\
\text{ensembles }(\tfrac12)\text{simpliciaux }\deleted{\text{\unsure{qco}}}\text{ (de Kan ?), groupes, anneaux etc}\end{gathered}\right.$
\margin{(idem)}
\unsure{Axiomatisation|Abstractisation} : catégorie $\frac12$ simpliciale

( $A$ catégorie, $C$ \deleted{\unsure{est}} catégorie (p. ex. ens. finis tot. ordonnés non vides \ldots)
la donnée d'un foncteur $\underset{\deleted{\unsure{\rm can}}}{\underline{\operatorname{Hom}}_{\cdot}}:A^o\times A\to\widehat C$
\inserted{$(A,B)\mapsto\underline{\operatorname{Hom}}_{\cdot}(A,B)$}
équivaut à celle d'un foncteur $A^o\times A\times C^o\mathbin{\unsure{-|\to}}\mathrm{Ens}$ ou encore
$C^o\times A\xrightarrow{\underline{\operatorname{Hom}}}\widehat A\quad(\sigma,B)\mapsto\underline{\operatorname{Hom}}(\sigma,B)$ ou d'un foncteur \inserted{$\sigma\mathbin{-}\varphi_\sigma$}
$C^o\longrightarrow\underline{\operatorname{Hom}}(A,\widehat A)\bigl[\cong\underline{\operatorname{Hom}}_{!}(\widehat A,\widehat A)\bigr]$ (le $!$ indique les
foncteurs qui commutent aux $\varinjlim$ quelc.) \deleted{soit} ou
encore un foncteur $\widehat{C}^{o}\to\underline{\operatorname{Hom}}(A,\widehat A)\simeq\underline{\operatorname{Hom}}_{!}(\widehat A,\widehat A)$
commutant aux $\varprojlim$ quelc, qu'on note encore $\sigma\to h_\sigma$.

Si $C$ a objet final $e$ (exemples dans le cas \unsure{spécifié|précisé})
alors la donnée d'un isom bifonctoriel
\[
\operatorname{Hom}_{\cdot}(A,B)(e)\simeq\operatorname{Hom}(A,B)\tag*{($\alpha$)}
\]
équivaut à la donnée d'un isom.
\[
\varphi_e=\mathrm{id}_{\widehat A}\tag*{($\beta$)}
\]

La donnée d'accouplements de trifoncteurs
\[
\underset{\deleted{\unsure{\rm can}}}{\operatorname{Hom}_{\cdot}}(A,B)\times\operatorname{Hom}_{\cdot}(B,C)\mathbin{-}\operatorname{Hom}_{\cdot}(A,C)\tag*{(a)}
\]
équivaut sauf erreur à celle d'hom\inserted{om}orphismes
bifonctoriels
\[
\varphi_{K\times L}\mathrel{\overset{\deleted{\unsure{\sim}}}{\leftarrow}}\varphi_K\circ\varphi_L\tag*{(b)}
\]
\margin{Quand \inserted{(b)} est-il un isomorphisme ? sss \textcircled{\scriptsize b$'$}
$\bigl\{$\diagram{\resizebox{\linewidth}{!}{\begin{tikzcd}[ampersand replacement=\&,row sep=small]
\operatorname{Hom}_{\cdot}(A,\underline{\operatorname{Hom}}(L,B))\arrow[d]\\
\underline{\operatorname{Hom}}(L,\operatorname{Hom}_{\cdot}(A,B))
\end{tikzcd}}}est un isom.}
(ou encore, d'hom fonctoriels
\[
\varphi_M\longleftarrow\varphi_{M'}\circ\varphi_N\tag*{(c)}
\]
)
et la \deleted{donnée d'un isom} \inserted{condition} d'associativité \deleted{\unsure{est}} pour
les accouplements entre les $\operatorname{Hom}_{\cdot}$ équivaut
sauf erreur à celle d'une \deleted{\unsure{isom}} compatibilité
d'associativité pour les données (b). Compatibilité
% archive sheet number: 25

\scanpage{27}
% transcription
entre donnée $\alpha$ et a), resp. $\beta$ et b) (ou c).
Dans le cas $C=$ ens. tot ordonnés $\ne\varnothing$, i.e. $\widehat C=$ ens. $\frac12$ simpl,
on a dans les $\operatorname{Hom}_{\cdot}(A,B)$ des relations d'homotopie.
Cela permet d'associer à $A$ des $n$-catégories pour tout $n$
(\unsure{Construction|construct} facile si les $\operatorname{Hom}_{\cdot}(A,B)$ sont des ens. $\frac12$ simpliciaux
de Kan \ldots)

\deleted{Hom} Définition de $\operatorname{Isom}_{\cdot}(A,B)$ et $\operatorname{Aut}_{\cdot}(A)$ (le premier
un groupe de $\widehat C$, le deuxième un pseudo-torseur sous
celui-ci). Notion d'homotopie de morphismes, d'isotopie
d'isomorphismes. Foncteurs \deleted{\unsure{de}} $\frac12$ simpliciaux de
catégories \inserted{à structure} semi-simpliciales \deleted{\unsure{Grp}} (Groupoïdes à structure
semi-simpliciale etc \ldots). On peut dire quand un
tel foncteur est une [$n$-]équivalence d'homotopie (
s'il induit \unsure{sur} les $\operatorname{Hom}_{\cdot}$ des ($n$)équivalences d'homotopie)
Sauf erreur, cela signifie aussi que le $n$-foncteur
entre les $n$-catégories associées est une $n$-équivalence.
\inserted{Cas des complexes d'une cat. abélienne. Lien avec ce qui précède via
Dold-Puppe.}

$2'$. \unsure{Tapis|Topis|Topos} de Quillen : interprétation des $n$-catégories
strictes comme structures $n$\inserted{fois}$\frac12$ simpliciales particulières.

3. $n$-catégories fibrées sur une catégorie, $n$-champs sur
un site un topos. \inserted{Ils forment une $(n+1)$-catégorie} Il faudrait voir

a) il y a une \inserted{$(n+1)$-}équivalence entre $n$-champs sur un
site, et $n$-champs sur le topos associé.

b) Toute $n$-catégorie fibrée sur un site admet un
$n$-champ enveloppant.

c) Existence d'images inverses \inserted{de $n$-champs} par morphismes
(de sites ou) de topos [ \uline{NB} La notion d'image directe
``est claire'' ; celle d'image inverse est censée s'en
déduire par \deleted{\unsure{une}} \inserted{\unsure{formule|formules} d'}adjonction
\[
\underline{\operatorname{Hom}}(A,f_*(B))\simeq\underline{\operatorname{Hom}}(f^*(A),B)
\]
[\unsure{$n$-équivalence|équivalence} \deleted{de} $n$-catégories \ldots].
% archive sheet number: 26

\scanpage{28}
% transcription
\deleted{\unsure{applications}}

4. \deleted{\unsure{les}} $n$-champs \inserted{sur $X$} de $Y$-profondeur $\geq i$ \quad $(0\leq i\leq n+2)$, $Y$
étant un fermé de $X$.

5. Formules de chgt de base en top. ordinaire \deleted{\unsure{ou|à}}
\margin{$n$-champs ind-finis constructibles etc \ldots}
ou en topologie étale \ldots\ (Il y en a au moins trois \deleted{\unsure{versions}} \inserted{types}
principaux \ldots)

6. \deleted{\unsure{Exemples}} \inserted{Théorèmes} du type de Lefschetz \deleted{\unsure{en}} en topologie
étale (schémas) ou dans l'analytique-complexe.

7. $n$-champs localement constants.
Un morphisme de topos \inserted{$f:X\to Y$} est une \inserted{$n$-}équivalence d'homotopie
sss pour tout $n$ champ loc ct. \inserted{$A$} sur $Y$,
il induit une $n$-équivalence
\[
\Gamma(Y,A)\mathrel{\overset{(n)}{\simeq}}\Gamma(X,f^*(A))
\]
On peut donner une version pour les types
d'homotopie ``complétés'' \unsure{pour|par} une catégorie de groupes,
à la Artin-Mazur ?

\deleted{\unsure{Ce}} [ Cela suggère qu'un ``type d'homotopie'' de topos
est incarné par la connaissance, pour tt $n$,
de la $(n+1)$-catégorie des $n$-champs loc. constants
dessus, avec les foncteurs de passage des \deleted{\unsure{champs}} $n$-champs aux $(n-1)$-champs sous-jacents et \unsure{inversem.|inversement}
Dans le cadre d'Artin-Mazur, il faudrait des
restrictions de constructibilité convenables sur les
$n$-champs avec lesquels on travaille \ldots

8. Structures algébriques sur des $n$-catégories et
$n$-champs : Gr-structure, structure de Picard,
de Picard-stricte, torseurs etc. Stabilité par images
directes et inverses de $n$-champs.

9. La $(n+1)$-catégorie des $n$-champs de Picard \inserted{stricts} sur $X$
est $(n+1)$-équivalente à la sous-catégorie pleine des
% archive sheet number: 27

\scanpage{29}
% transcription
la catégorie dérivée $D(\operatorname{Ab}(X))$ formée des complexes
\deleted{\unsure{à cohomologie}} \deleted{qui sont} tels que $\underline H^i(K^{\unsure{\cdot|'}})=0$ si $i\notin[0,n]$. Les
opérations image inverse et image directe \deleted{sont} sur les
$n$-champs de Picard strict correspondent aux
opérations analogues dans les catégories dérivées (sauf
qu'il faut tronquer, dans le cas de l'image
directe). Les $\underline{\operatorname{Hom}}$ de cat. de Picard \unsure{stricts|strictes}
correspondent aux $\mathrm R\underline{\operatorname{Hom}}$ tronqués \inserted{sur $[0,n]$} dans les
catégories dérivées.

\uline{NB} Il y aurait lieu d'interpréter de façon
géométrique les objets tronqués de $D(\operatorname{Mod}(O))$,
$O$ un Anneau sur $X$, en termes de cat. de
Picard strictes où $O$ opère, et d'interpréter géométriquement les opérations $\mathrm Rf_*$, \unsure{de|et} $\mathrm Lf^*$ et $\mathrm R\underline{\operatorname{Hom}}$.
Interprétation des $\overset{(\mathrm L)}{\otimes}_{\unsure{k|Z}}$ en termes d'accouplements
\inserted{``universels''} de cat. de Picard ?

10. Etude des $n$-Gr champs. \deleted{\unsure{E}} J'ai regardé le
cas $n=1$, $X$ quelconque\inserted{, cf lettre à Deligne}. Mais \unsure{même} si $X$ est le
topos ponctuel, je ne sais quoi dire pour le
cas $n=2$ p.ex. Il faudrait examiner le cas
de la $n$-\inserted{Gr-}catégorie associée à un groupe $\frac12$ simplicial (p.ex. via un groupe topologique), cf \S~2, et voir
la $n$-catégorie des homomorphismes de \deleted{\unsure{groupes}} \inserted{l'objet}
$\frac12$ simplicial unité dans ledit groupe $\frac12$ simplicial.

11. $n$-\deleted{catégorie}\inserted{champ} de Picard \inserted{(pas strict !)} enveloppant d'un $(n)$-$\otimes$ champ $A$
ACU. Ses invariants $\pi_i$ sont les faisceaux $\underline{K}_i(A)$.
% archive sheet number: 28

\scanpage{30}
% transcription
\uline{Question} Que peut-on dire de la structure des $n$-champs de Picard pas stricts ? On voit pour $n=1$
(cf correspondance avec Breen et \uline{Deligne}), mais
j'ignore pour $n\geq2$ \ldots
% archive sheet number: 29
